\documentclass[12pt]{article}
\usepackage[utf8]{inputenc}
\usepackage[T1]{fontenc}
\usepackage{hyperref}
\title{Wom Language Resources for K-12 Curriculum Development}
\author{Compiled by Claude Sonnet 5.5 (claude-sonnet-5-5) for Andr\'as Kornai}
\date{}
\begin{document}
\maketitle

\section{Introduction}
Wom (ISO 639-3 \texttt{wom}; also called Perema, Pere and Pereba) is a Leko language of Adamawa State, Nigeria. Most sources place it in the Fufore area. The searches found no evidence that it is spoken in Bauchi State. Speaker estimates differ: about 5,000 (Wikipedia, ASJP), 8,580 (Global Recordings Network), and 12,000 in Nigeria (Joshua Project). The community includes both Christians and Muslims, and speakers are reported to be literate in Hausa and English. Very little material exists.

\section{Primary Educational Resources}
No primers, textbooks or school programmes in Wom were found.

\section{Linguistic and Reference Materials}
\begin{itemize}
\item Wikipedia, Perema language: \url{https://en.wikipedia.org/wiki/Perema_language}
\item Wikipedia, Leko languages: \url{https://en.wikipedia.org/wiki/Leko_languages}
\item Numerals of Wom, Max Planck lingweb: \url{https://lingweb.eva.mpg.de/channumerals/Wom.htm}
\item The 2017 CALL AdaGram survey report describes a wordlist of about 500 words with audio and about 140 phrases with audio, and notes grammar and documentation work by Jakob Lesage. I did not see where the wordlist itself is published. \url{https://halshs.archives-ouvertes.fr/halshs-01676457v1/file/2017_CALL_AdaGram_Survey_Results.pdf}
\item Glottolog languoid page: \url{https://glottolog.org/resource/languoid/id/womn1235.bigmap.html}
\end{itemize}

\section{Religious and Community Resources}
\begin{itemize}
\item Joshua Project language page, Scripture status ``Translation Started''; no audio Bible: \url{https://joshuaproject.net/index.php/languages/wom}
\item Joshua Project, Wom in Nigeria: \url{https://www.joshuaproject.net/index.php/people_groups/15415/NI}
\end{itemize}

\section{Digital and Audio Resources}
\begin{itemize}
\item Global Recordings Network, Wom [Nigeria]: evangelism and Bible-teaching audio (MP3 set and slideshow): \url{https://globalrecordings.net/en/language/wom}
\end{itemize}

\section{Curriculum Development Context}
Wom is not known to be taught in schools. Documentation is limited to survey wordlists, numerals and a documentation project, and no orthography was found. A curriculum would start from community-led orthography work and recordings, with Hausa and English as the school languages of wider communication.
\end{document}